\begin{abstract}
\noindent
A previous FO-LQI $\to$ FO-PIDF optimisation campaign for a non-ideal 48~V
Buck converter spent more than 1800 evaluations without reaching robust
feasibility, and attributed this failure to clipping of the proportional gain
at a parameterisation bound. We propose a \emph{bound-saturation audit} that
separates structural from pathological saturation, and apply it to a
reconstructed chain whose projection reproduces the exact identity
proposition to within \ChiffreErreurPropUn. Over \ChiffreAuditN{} samples of
the decision box, $\log_{10}k_p$ never reaches its bound: clipping does not
explain the failure. A preregistered confirmatory campaign (hash-frozen
protocol, three arms, six metaheuristics, thirty paired seeds, 540 cells) on a
switched simulator that reproduces discontinuous conduction yields \emph{no}
robustly feasible controller: 0/30 seeds per method and per arm, Wilson
interval $[0;\ChiffreWilsonHaut]$. The parameterisation specified by the
brief does not even reach nominal feasibility; freeing the setpoint weight $b$
makes it reachable in \ChiffreNomFaisCinqB{} of 180 cells, without
robustness. A physical diagnosis identifies a controller-independent
obstruction: at minimum duty ratio, the converter in discontinuous conduction
settles below the minimum-voltage threshold in \ChiffrePlancherN{} of the
\ChiffrePlancherTot{} uncertainty cases. We further show that controllers
feasible with respect to all six constraints can sustain a limit cycle that
these constraints do not detect. These results concern a simulated model and
a finite set of cases; they are neither a robustness certificate nor a proof
of infeasibility.
\end{abstract}

\noindent\textbf{Keywords:} Buck converter; fractional-order control; FO-LQI;
FO-PIDF; metaheuristics; parameterisation audit; discontinuous conduction;
reproducibility.
